% Bài tự luận (chặng Tự luận, chat cùng Cherry), lớp 9 chương 5. Dựa trên docs/mau/bai-mau-chuong.tex, thêm \huong, \goiy, \kiemtra sau mỗi \buocgiai; bỏ sơ đồ (thầy chọn không dùng sơ đồ đi lên).
% Mẫu bài mẫu theo chương: tên bài mẫu, dòng saubai ghi mã bài (mở cho học sinh khi qua bài đó; bỏ dòng này thì mở khi qua hết chương), rồi khối giaimau.
% Đã tự kiểm tra: tọa độ hình tính bằng Python (tiếp điểm, vuông góc, trung điểm, giao điểm đều đúng, sai số do làm tròn căn 3); OH.OA = R^2 chứng minh bằng tam giác đồng dạng (lớp 8), không dùng hệ thức lượng.
\chude{Hai tiếp tuyến cắt nhau}
\saubai{k9-b16}
\begin{giaimau}
\de{Cho đường tròn $(O;R)$ và điểm $A$ nằm ngoài đường tròn. Kẻ hai tiếp tuyến $AB$, $AC$ với đường tròn ($B$, $C$ là các tiếp điểm). Gọi $H$ là giao điểm của $OA$ và $BC$. Kẻ đường kính $BD$ của đường tròn $(O)$.\\
1) Chứng minh bốn điểm $A$, $B$, $O$, $C$ cùng thuộc một đường tròn và $OH\cdot OA=R^2$.\\
2) Chứng minh $CD\parallel OA$.\\
3) Kẻ $CK\perp BD$ ($K$ thuộc $BD$). Đoạn thẳng $AD$ cắt $CK$ tại $I$. Chứng minh $I$ là trung điểm của $CK$.}
\begin{tikzpicture}[scale=1, line join=round]
\tkzSetUpLabel[font=\footnotesize]
\tkzSetUpLine[line width=.8pt, color=black]
\tkzDefPoint(0,0){O}
\tkzDefPoint(4,0){A}
\tkzDefPoint(1,1.7321){B}
\tkzDefPoint(1,-1.7321){C}
\tkzDefPoint(-1,-1.7321){D}
\tkzDefPoint(1,0){H}
\tkzDefPoint(-0.5,-0.8660){K}
\tkzDefPoint(0.25,-1.2990){I}
\tkzDefPoint(7,-1.7321){M}
\tkzDefPoint(2,0){J}
\tkzDrawCircle[line width=.8pt, color=black](O,B)
\tkzDrawSegments(A,B A,C B,C O,A O,B O,C)
\tkzMarkRightAngles[size=.18](O,B,A O,C,A)
\phu{2}{3}{\draw[dashed] (J) circle (2); \draw[dashed] (J)--(B) (J)--(C); \draw[fill=black] (J) circle (.04) +(90:.3) node[font=\footnotesize]{$J$};}
\hien{4}{\tkzMarkRightAngles[size=.18](B,H,O)}
\hien{6}{\tkzDrawSegments(B,D C,D) \draw[fill=black] (D) circle (.04) +(225:.3) node[font=\footnotesize]{$D$};}
\hien{8}{\tkzDrawSegments(C,K A,D) \tkzMarkRightAngles[size=.18](C,K,D) \foreach \x/\g in {K/170, I/-110} \draw[fill=black] (\x) circle (.04) +(\g:.3) node[font=\footnotesize]{$\x$};}
\hien{8}{\tkzDrawSegments[dashed](A,M C,M) \draw[fill=black] (M) circle (.04) +(-20:.3) node[font=\footnotesize]{$M$};}
\hien{13}{\tkzMarkSegments[mark=|](C,I I,K)}
\sang{1}{\draw (A)--(B)--(O)--(C)--cycle;}
\sang{2}{\draw (J)--(B) (J)--(C) (O)--(A);}
\sang{3}{\draw (J) circle (2);}
\sang{4}{\draw (A)--(B) (A)--(C) (O)--(B) (O)--(C) (B)--(C);}
\sang{5}{\draw (O)--(B)--(A)--cycle (B)--(H);}
\sang{6}{\draw (B)--(C)--(D)--cycle (O)--(C);}
\sang{7}{\draw (C)--(D) (O)--(A) (B)--(C);}
\sang{8}{\draw[dashed] (A)--(M)--(C);}
\sang{9}{\draw (B)--(M)--(D)--cycle (O)--(A);}
\sang{10}{\draw (C)--(K) (B)--(M) (B)--(D);}
\sang{11}{\draw (D)--(B)--(A)--cycle (I)--(K);}
\sang{12}{\draw (D)--(M)--(A)--cycle (I)--(C);}
\sang{13}{\draw (C)--(K);}
\foreach \x/\g in {O/150, A/60, B/90, C/-80, H/60}
  \draw[color=black, fill=black] (\x) circle (.04) +(\g:.3) node[font=\footnotesize]{$\x$};
\end{tikzpicture}
\cau{Câu 1}{Chứng minh bốn điểm $A$, $B$, $O$, $C$ cùng thuộc một đường tròn và $OH\cdot OA=R^2$.}
\doan{Muốn chứng minh bốn điểm cùng thuộc một đường tròn, em nên tìm gì?}{Một điểm cách đều cả bốn điểm $A$, $B$, $O$, $C$}{Chứng minh $ABOC$ là hình vuông}{Điểm cách đều bốn điểm chính là tâm của đường tròn đi qua chúng.}{$ABOC$ nói chung không là hình vuông ($OB=R$ còn $AB$ thường khác $R$). Hãy tìm một điểm cách đều bốn điểm.}{Chứng minh $OA=BC$}{Đẳng thức này nói chung không đúng, và cũng không chỉ ra được đường tròn nào.}
\buocgiai{Tìm hai góc vuông có sẵn}{$AB$, $AC$ là tiếp tuyến tại $B$, $C$ nên $AB\perp OB$, $AC\perp OC$.}{Tiếp tuyến vuông góc với bán kính tại tiếp điểm.}
\huong{Đề cho $AB$, $AC$ là tiếp tuyến. Tiếp tuyến luôn vuông góc với bán kính đi qua tiếp điểm. Cậu viết ra hai quan hệ vuông góc nhé.}
\goiy{Tiếp tuyến tại $B$ vuông góc với đoạn nào cũng đi qua $B$?}
\goiy{Đó là bán kính $OB$. Ở điểm $C$ là bán kính $OC$.}
\kiemtra{Góc nào sau đây chắc chắn bằng $90^\circ$?}{$\widehat{ABO}$}{$\widehat{BAC}$}{Vì $AB\perp OB$ tại tiếp điểm $B$.}{Góc $BAC$ tùy vị trí điểm $A$, đề không cho nó vuông.}{$\widehat{BOC}$}{Góc $BOC$ cũng tùy vị trí điểm $A$. Hãy nhìn chỗ tiếp tuyến gặp bán kính.}
\buocgiai{Lấy trung điểm của cạnh huyền chung}{Gọi $J$ là trung điểm của $OA$.\\ $\triangle OBA$ vuông tại $B$ nên $JB=\dfrac{OA}{2}$.\\ $\triangle OCA$ vuông tại $C$ nên $JC=\dfrac{OA}{2}$.}{Trong tam giác vuông, đường trung tuyến ứng với cạnh huyền bằng nửa cạnh huyền.}
\huong{Hai tam giác vuông $OBA$ và $OCA$ có chung cạnh huyền $OA$. Cậu gọi $J$ là trung điểm của $OA$, rồi nhớ lại tính chất đường trung tuyến ứng với cạnh huyền.}
\goiy{Trong tam giác vuông, đường trung tuyến ứng với cạnh huyền bằng nửa cạnh huyền.}
\kiemtra{Trong tam giác vuông $OBA$, đoạn $JB$ bằng bao nhiêu?}{$\dfrac{OA}{2}$}{$\dfrac{AB}{2}$}{$BJ$ là trung tuyến ứng với cạnh huyền $OA$.}{$BJ$ là trung tuyến ứng với cạnh huyền $OA$, nên phải so với $OA$ chứ không phải $AB$.}{$OB$}{$JB=OB$ chỉ đúng trong trường hợp đặc biệt, không dùng được.}
\buocgiai{Kết luận bốn điểm}{$JA=JO=JB=JC=\dfrac{OA}{2}$.\\ Vậy $A$, $B$, $O$, $C$ cùng thuộc đường tròn tâm $J$, đường kính $OA$.}{Bốn điểm cách đều điểm $J$ thì cùng nằm trên đường tròn tâm $J$.}
\huong{Giờ cậu so bốn đoạn $JA$, $JO$, $JB$, $JC$ xem.}
\goiy{$J$ là trung điểm của $OA$ nên $JA=JO=\dfrac{OA}{2}$.}
\kiemtra{Bốn điểm $A$, $B$, $O$, $C$ cùng thuộc đường tròn nào?}{Tâm $J$, đường kính $OA$}{Tâm $O$, bán kính $R$}{Cả bốn điểm cách $J$ một khoảng bằng $\dfrac{OA}{2}$.}{$A$ nằm ngoài $(O)$ nên không thuộc $(O)$.}{Tâm $H$, bán kính $HB$}{Nói chung $HA\ne HB$, nên $A$ không nằm trên đường tròn này.}
\doan{Để có $OH\cdot OA=R^2$, em cần biết thêm điều gì về $BH$?}{$BH\perp OA$, tức $BH$ là đường cao của $\triangle OBA$}{$BH=HC$}{Có $BH\perp OA$ thì $\triangle OHB$ và $\triangle OBA$ cùng vuông, có chung góc $O$, nên đồng dạng.}{$BH=HC$ đúng, nhưng tự nó chưa cho tích $OH\cdot OA$. Cần một quan hệ vuông góc để có hai tam giác vuông đồng dạng.}{$H$ là trung điểm của $OA$}{Nói chung $H$ không là trung điểm của $OA$ (trên hình $H$ gần $O$ hơn).}
\buocgiai{Chứng minh $OA\perp BC$}{$AB=AC$ (tính chất hai tiếp tuyến cắt nhau), $OB=OC=R$.\\ Nên $OA$ là đường trung trực của $BC$, suy ra $OA\perp BC$ tại $H$.}{$A$ và $O$ cùng cách đều $B$, $C$ nên cùng nằm trên đường trung trực của $BC$.}
\huong{Điểm $A$ có cách đều $B$ và $C$ không? Còn điểm $O$ thì sao? Hai điểm cùng cách đều $B$, $C$ thì nằm trên đường gì của đoạn $BC$?}
\goiy{$AB=AC$ theo tính chất hai tiếp tuyến cắt nhau, còn $OB=OC=R$.}
\goiy{Hai điểm cùng cách đều hai đầu đoạn $BC$ thì cùng nằm trên đường trung trực của $BC$.}
\kiemtra{Vì sao $AB=AC$?}{Tính chất hai tiếp tuyến cắt nhau}{Vì $\triangle ABC$ vuông}{Hai đoạn tiếp tuyến kẻ từ một điểm nằm ngoài đường tròn thì bằng nhau.}{Đề không cho $\triangle ABC$ vuông.}{Vì $H$ là trung điểm của $BC$}{Đó là điều suy ra sau, chưa dùng được ở bước này.}
\buocgiai{Dùng tam giác đồng dạng}{$\triangle OHB$ và $\triangle OBA$ có $\widehat{OHB}=\widehat{OBA}=90^\circ$, chung góc $O$.\\ Nên $\triangle OHB\backsim\triangle OBA$ (g.g), suy ra $\dfrac{OH}{OB}=\dfrac{OB}{OA}$.\\ Do đó $OH\cdot OA=OB^2=R^2$.}{Hai tam giác vuông có một góc nhọn bằng nhau thì đồng dạng; tỉ số các cạnh tương ứng bằng nhau.}
\huong{Hai tam giác vuông $OHB$ và $OBA$ có chung góc $O$. Cậu viết cặp tam giác đồng dạng, rồi lập tỉ số các cạnh tương ứng.}
\goiy{$\widehat{OHB}=\widehat{OBA}=90^\circ$ và chung góc $O$ nên $\triangle OHB\backsim\triangle OBA$ (g.g).}
\goiy{Các đỉnh tương ứng: $O$ với $O$, $H$ với $B$, $B$ với $A$.}
\kiemtra{Từ $\triangle OHB\backsim\triangle OBA$, tỉ số nào đúng?}{$\dfrac{OH}{OB}=\dfrac{OB}{OA}$}{$\dfrac{OH}{OA}=\dfrac{OB}{HB}$}{Cạnh $OH$ ứng với $OB$, cạnh $OB$ ứng với $OA$. Nhân chéo được $OH\cdot OA=OB^2$.}{Ghép sai cặp cạnh. Các đỉnh tương ứng là $O$ với $O$, $H$ với $B$, $B$ với $A$.}{$\dfrac{OH}{HB}=\dfrac{OA}{OB}$}{Ghép sai: $\dfrac{OH}{HB}$ phải bằng $\dfrac{OB}{BA}$.}
\cau{Câu 2}{Chứng minh $CD\parallel OA$.}
\doan{Hai đường thẳng $CD$ và $OA$ cùng có quan hệ gì với một đường thẳng khác?}{Cùng vuông góc với $BC$}{Cùng đi qua $O$}{$OA\perp BC$ đã có ở câu 1, chỉ còn chứng minh $CD\perp BC$.}{$CD$ không đi qua $O$.}{Cùng song song với $AB$}{$OA$ cắt $AB$ tại $A$ nên không song song với $AB$.}
\buocgiai{Chứng minh $CD\perp BC$}{$O$ là trung điểm của $BD$ nên $CO$ là trung tuyến của $\triangle BCD$.\\ $CO=R=\dfrac{BD}{2}$ nên $\triangle BCD$ vuông tại $C$, tức $CD\perp BC$.}{Tam giác có đường trung tuyến bằng nửa cạnh tương ứng thì vuông tại đỉnh đó.}
\huong{Nhìn tam giác $BCD$: $O$ là trung điểm của $BD$, và $OC=R$. Cậu so $OC$ với $BD$ nhé.}
\goiy{$BD=2R$ nên $OC=\dfrac{BD}{2}$.}
\kiemtra{$CO=\dfrac{BD}{2}$ cho ta biết $\triangle BCD$ vuông tại đâu?}{Tại $C$}{Tại $B$}{Góc vuông nằm ở đỉnh đối diện với cạnh $BD$.}{Góc vuông nằm ở đỉnh đối diện với cạnh $BD$, tức là đỉnh $C$.}{Tại $O$}{$O$ không phải đỉnh của tam giác $BCD$.}
\buocgiai{Kết luận}{$CD\perp BC$ và $OA\perp BC$ (câu 1).\\ Suy ra $CD\parallel OA$.}{Hai đường thẳng cùng vuông góc với một đường thẳng thì song song với nhau.}
\huong{Cậu đã có hai đường thẳng cùng vuông góc với $BC$. Viết kết luận thôi!}
\cau{Câu 3}{Kẻ $CK\perp BD$ ($K$ thuộc $BD$). Đoạn thẳng $AD$ cắt $CK$ tại $I$. Chứng minh $I$ là trung điểm của $CK$.}
\doan{Câu 2 cho $CD\parallel OA$, nhưng trên hình chưa dùng được ngay. Em nên vẽ thêm gì?}{Kéo dài $DC$ cắt đường thẳng $AB$ tại $M$}{Kẻ qua $I$ đường thẳng song song với $BD$}{Khi đó $OA\parallel DM$ nằm trong $\triangle BDM$, dùng được định lí Thales.}{Đường này không dùng tới $CD\parallel OA$, khó dẫn tới $IK=IC$.}{Nối $O$ với $I$}{$OI$ không tạo ra tam giác nào chứa cặp song song $CD\parallel OA$.}
\buocgiai{Tạo điểm $M$}{Gọi $M$ là giao điểm của đường thẳng $DC$ và đường thẳng $AB$.}{Vẽ thêm điểm để biến giả thiết $CD\parallel OA$ thành cặp song song nằm trong một tam giác.}
\huong{Viết câu “Gọi $M$ là giao điểm của đường thẳng $DC$ và đường thẳng $AB$” trước, rồi mới dùng điểm $M$.}
\buocgiai{Chứng minh $AB=AM$}{Trong $\triangle BDM$: $O$ là trung điểm của $BD$, $OA\parallel DM$ (câu 2), $A\in BM$.\\ Theo định lí Thales: $\dfrac{BA}{AM}=\dfrac{BO}{OD}=1$, nên $AB=AM$.}{Đường thẳng đi qua trung điểm một cạnh và song song với cạnh thứ hai thì đi qua trung điểm cạnh thứ ba.}
\huong{Trong $\triangle BDM$: $O$ là trung điểm của $BD$ và $OA\parallel DM$. Đường thẳng đi qua trung điểm một cạnh, song song với cạnh thứ hai thì cắt cạnh thứ ba ở đâu?}
\goiy{Nó đi qua trung điểm của cạnh thứ ba, tức là $A$ là trung điểm của $BM$.}
\kiemtra{Điểm $A$ là gì của đoạn $BM$?}{Trung điểm của $BM$}{Chân đường cao kẻ từ $D$}{Nên $AB=AM$.}{Không có quan hệ vuông góc nào cho $DA\perp BM$.}{Điểm chia $BM$ theo tỉ số $1:2$}{Tỉ số là $\dfrac{BA}{AM}=\dfrac{BO}{OD}=1$, tức là $1:1$.}
\buocgiai{Chứng minh $CK\parallel BM$}{$CK\perp BD$ (giả thiết) và $BM\perp BD$ (vì $AB$ là tiếp tuyến tại $B$).\\ Nên $CK\parallel BM$, tức là $IK\parallel AB$ và $IC\parallel AM$.}{Hai đường thẳng cùng vuông góc với $BD$ thì song song.}
\huong{$CK\perp BD$ theo đề. Còn $BM$, tức là tiếp tuyến tại $B$, có vuông góc với $BD$ không?}
\goiy{$BD$ là đường kính nên đi qua tiếp điểm $B$.}
\kiemtra{Vì sao $BM\perp BD$?}{Vì $AB$ là tiếp tuyến tại $B$ và $BD$ là đường kính}{Vì $\triangle BDM$ cân}{Tiếp tuyến vuông góc với bán kính tại tiếp điểm.}{Không có dữ kiện nào cho $\triangle BDM$ cân.}{Vì $CD\parallel OA$}{$CD\parallel OA$ không nói gì về góc giữa $BM$ và $BD$.}
\doan{Có $IK\parallel AB$ và $IC\parallel AM$. Hai tam giác nào cho cùng tỉ số $\dfrac{DI}{DA}$?}{$\triangle DBA$ và $\triangle DMA$}{$\triangle DKI$ và $\triangle DCI$}{Hai tam giác chung đỉnh $D$ và chung cạnh $DA$ chứa điểm $I$.}{Hai tam giác này không chứa $AB$, $AM$ nên không nối được với $AB=AM$.}{$\triangle OAB$ và $\triangle OAC$}{Hai tam giác này không chứa điểm $I$.}
\buocgiai{Thales trong $\triangle DBA$}{$IK\parallel AB$ với $I\in DA$, $K\in DB$:\\ $\dfrac{IK}{AB}=\dfrac{DI}{DA}$. (1)}{Đường thẳng song song với một cạnh tạo ra tam giác mới có ba cạnh tương ứng tỉ lệ.}
\huong{Trong $\triangle DBA$ có $IK\parallel AB$. Cậu viết tỉ số, tính từ đỉnh chung $D$.}
\goiy{Tam giác nhỏ là $DKI$, tam giác lớn là $DBA$.}
\kiemtra{Trong $\triangle DBA$ với $IK\parallel AB$, tỉ số nào đúng?}{$\dfrac{IK}{AB}=\dfrac{DI}{DA}$}{$\dfrac{IK}{AB}=\dfrac{IA}{DA}$}{Đúng hệ quả định lí Thales.}{Tam giác nhỏ là $DKI$ có đỉnh $D$, nên phải là $\dfrac{DI}{DA}$.}{$\dfrac{IK}{AB}=\dfrac{DK}{KB}$}{Tỉ số đúng phải là $\dfrac{DK}{DB}$, không phải $\dfrac{DK}{KB}$.}
\buocgiai{Thales trong $\triangle DMA$}{$IC\parallel AM$ với $I\in DA$, $C\in DM$:\\ $\dfrac{IC}{AM}=\dfrac{DI}{DA}$. (2)}{Cùng hệ quả Thales, lần này trong tam giác $DMA$.}
\huong{Làm tương tự trong $\triangle DMA$ với $IC\parallel AM$.}
\goiy{Tam giác nhỏ là $DCI$, tam giác lớn là $DMA$.}
\kiemtra{Trong $\triangle DMA$ với $IC\parallel AM$, $\dfrac{IC}{AM}$ bằng tỉ số nào?}{$\dfrac{DI}{DA}$}{$\dfrac{DC}{CM}$}{Cùng hệ quả định lí Thales, tỉ số tính từ $D$.}{$\dfrac{DC}{CM}$ là tỉ số hai đoạn chia; tỉ số đúng là $\dfrac{DC}{DM}$.}{$\dfrac{IA}{DA}$}{Tam giác nhỏ $DCI$ có đỉnh $D$, nên phải là $\dfrac{DI}{DA}$.}
\buocgiai{Kết luận}{Từ (1) và (2): $\dfrac{IK}{AB}=\dfrac{IC}{AM}$.\\ Mà $AB=AM$ nên $IK=IC$.\\ Vậy $I$ là trung điểm của $CK$.}{Hai phân số bằng nhau có mẫu bằng nhau thì tử bằng nhau.}
\huong{Ghép $(1)$ với $(2)$, rồi dùng $AB=AM$.}
\kiemtra{Từ $(1)$, $(2)$ và $AB=AM$, ta suy ra được gì?}{$IK=IC$}{$IK=AB$}{Vậy $I$ là trung điểm của $CK$.}{Hai tỉ số bằng nhau chứ $IK$ không bằng $AB$.}{$DI=DA$}{$I$ nằm giữa $D$ và $A$ nên $DI<DA$.}
\nhanxet{Ý phân loại: phải tạo thêm điểm $M$ để dùng được giả thiết $CD\parallel OA$ của câu 2.\\ Và phải nhìn ra hai tam giác $DBA$, $DMA$ có chung tỉ số $\dfrac{DI}{DA}$.}
\end{giaimau}
